%%%%%%%%%%%%%%%%%%%%%%%%%%%%%%%%%%%%%%%%%%%%%%%%%%%%%%%%%%%%%%%%%%
%% Toto je inkludovaný LaTeXový soubor `x_diag.tex'.            %%
%% This is the `x_diag.tex' included LaTeX file.                %%
%% ------------------------------------------------------------ %%
%% Verze / Version 2022-05-23                                   %%
%% Vyvíjí & udržuje / Developed & maintained by                 %%
%%                           <stanislav.hledik@physics.slu.cz>  %%
%%%%%%%%%%%%%%%%%%%%%%%%%%%%%%%%%%%%%%%%%%%%%%%%%%%%%%%%%%%%%%%%%%

\chapter{Diagnostika s~obrazovým výstupem}\label{princtwo}%################

V~oblasti lékařské diagnostiky\INDEX{diagnostika!lékařská} se v~posledních
letech uplatňují nové metody zobrazování, jež by nemohly vzniknout bez
počítačové podpory při transformaci primárních dat produkovaných fyzikálními
detektory na diagnosticky upotřebitelné obrazové výstupy. Jedná se
ultrazvukovou sonografii (US), počítačovou tomografii (CT), magnetickou
rezonanci (MRI, fMRI), radioizotopové scintigrafické metody (planární
scintigrafie, SPECT, PET).  Všechny uvedené metody jsou v~současné medicínské
praxi úspěšně používány díky expandujícímu rozvoji, který v~poslední době
nastal hlavně v~oblasti detektorů záření a informačních technologií při
zpracování velkých objemů dat. V~oblasti diagnostiky je proto nutná úzká
spolupráce mezi medicínou, fyzikou a informatikou.

Z~pochopitelných důvodů je zde prostor pouze pro fyzikální a informatické
základy uvedených metod. Další poučení může čtenář získat např.
z~monografií~\cite{Dow-Ken-Joh:2006:PhysDiagImag:,%
  Roz:2006:ElePristLekar:,%
  Nav-etal:2005:MedBiofyz:}. V~této kapitole se soustředíme na popis
ultrazvukové sonografie a zobrazování pomocí (jaderné) magnetické rezonance.

\section{Ultrazvuková sonografie}\label{sono}%%%%%%%%%%%%%%%%%%%%%%%%%%%%%%

Zvukové vlny vysoké frekvence ($2.5$--$40\,\hz[M]$)\footnote{Pro srovnání,
  horní mez slyšitelnosti je pro lidské ucho okolo $20\,\hz[k]$, u~dětí až
  $40\,\hz[k]$.} mohou poskytovat 2D obrazové informace o~tkáních \emph{in
  vivo}. Dají se využít také pro funkční zobrazování v~živých tkáních, např.
na základě Dopplerova jevu. Zobrazování pomocí ultrazvuku (\emph{sonografie})
je podobné nejběžnější zobrazovací metodě pomocí X-záření\INDEX{záření!X}
v~tom, že využívá fyzikálního agens schopného prostupovat (pokud možno bez
poškození) tkáněmi, liší se však v~několika základních aspektech:
\begin{itemize}
\item Ultrazvuk, jakožto mechanické vlnění, potřebuje látkové prostředí,
  ve kterém se může šířit.
\item Na rozdíl od X-záření,\INDEX{záření!X} jež je zářením
  elektromagnetickým,\INDEX{záření!elektromagnetické} je ultrazvuk podélným
  vlněním, a~jeho vlnová délka je podstatně větší. Pro frekvenci $2\,\hz[M]$
  je vlnová délka v~měkké tkáni $0.75\,\mathrm{mm}$, pro frekvenci
  $10\,\hz[M]$ je vlnová délka $0.15\,\mathrm{mm}$. Pro srovnání,
  charakteristické vlnové délky X-záření\INDEX{záření!X} při urychlovacím
  napětí okolo $100\,\mathrm{kV}$ jsou $10^{7}\times$ větší.
\item Pro zobrazování se namísto prošlého vlnění jako u~rentgenového
  zobrazování využívá odraženého vlnění.
\end{itemize}

\begin{figure}[t]
\begin{center}

\includegraphics[width=.6\linewidth]{echoloc}
\end{center}
\caption[K~echolokačnímu principu ultrazvukové sonografie]%
{\label{echoloc}K~echolokačnímu principu ultrazvukové sonografie. Ultrazvukový
  puls odražený od vzdálenějšího rozhraní potřebuje delší dobu pro návrat echa
  a vyvolání odezvy v~detektoru. Z~časové prodlevy mezi vysláním pulsu a
  detekcí echa lze lokalizovat rozhraní. Převzato
  z~\cite{Dow-Ken-Joh:2006:PhysDiagImag:}.}
\end{figure}

Ultrazvukový signál je vytvářen elektrickou stimulací krystalického materiálu.
Jeho mechanické oscilace se přes vrstvu kontaktního gelu přenášejí do tkáně.
Vzduch a různé druhy tkání (měkké tkáně, tuk, krev, mozek, svaly, kosti,
apod.) se liší mechanickými vlastnostmi ovlivňujícími šíření ultrazvukových
vln. Na rozhraní dvou různých tkání dochází k~více či méně silnému odrazu
ultrazvukové vlny. Zdroj ultrazvukového vlnění (transducer) vysílá posloupnost
krátkých ultrazvukových pulsů a po vyslání každého pulsu se přepne do
\uv{naslouchacícho} režimu, v~němž zaznamenává odražené pulsy (echa)
přicházející z~nitra vyšetřovaného organismu. Na základě časové prodlevy mezi
vysláním pulsu a detekcí echa, intenzity echa a směrů, z~nichž echa
přicházejí, se konstruuje obraz. Tento echolokační princip je schematicky
zachycen na obrázku~\ref{echoloc}.\footnote{Princip ultrazvukové echolokace
  není v~přírodě neznámý~-- ultrazvukových vln o~frekvencích
  $50$--$200\,\hz[k]$ používají k~echolokaci netopýři. Vlastní medicínská
  aplikace ultrazvuku se vyvinula z~vojenské technologie~-- sonaru (\uv{sound
    navigation and ranging}) pro lokalizaci nepřátelských
  ponorek~\cite{Zaz:2007::NaziKost}. Maximální frekvence zvuku leží okolo
  $600\,\hz[M]$~\cite{Dow-Ken-Joh:2006:PhysDiagImag:}.}

\subsection{Základní pojmy ultrazvukové akustiky}\label{accoustic}%////////

V~této sekci připomeneme základní pojmy z~akustiky, jež jsou nezbytné pro
pochopení ultrazvukové sonografie.

V~plynném či kapalném prostředí se zvuk šíří ve formě \emph{podélných}
(longitudálních) vln, v~pevném skupenství je možná i~forma \emph{příčných}
(transverzálních) vln. Vzhledem k~oblasti aplikací ultrazvuku v~medicíně~--
vyšetření měkkých tkání, jež se svými vlastnostmi blíží spíše kapalině než
pevnému skupenství~-- nás zde zajímají jen podélné vlny. Podélné vlny si
můžeme představit takto: zařídíme-li nějakými vhodnými okrajovými podmínkami
mechanické oscilace na povrchu tkáně, energie těchto oscilací se šíří do tkání
postupným přenášením pohybu na sousedící molekuly tak, že výchylka molekul je
rovnoběžná se směrem šíření rozruchu.\footnote{Molekuly kmitají podél směru
  šíření, odtud název podélné vlnění.} Ačkoli samotné molekuly konají v~daném,
zafixovaném místě jen malý oscilační pohyb okolo rovnovážné polohy, samotný
oscilační rozruch, zahrnující kromě změny polohy molekul také malou změnu
rychlosti, hustoty a tlaku, se šíří \emph{fázovou rychlostí} $c$ zavisející na
mechanických vlastnostech tkáně~-- takový pohyb nazýváme vlnovým. Kinematicky
je vlnový pohyb popsán \emph{frekvencí} $f$ (počet cyklů detekovaných
pozorovatelem na fixním místě za jednu sekundu), případně periodou $T=1/f$,
vlnovou délkou
\begin{equation}
  \lambda = \frac{c}{f}                                     \label{sonwlen}
\end{equation}
a amplitudou $A$.  Názorná představa podélné zvukové vlny je znázorněna na
obrázku~\ref{longwav}. Rychlosti šíření pro některá medicínsky důležitá
prostředí jsou uvedeny v~tabulce~\ref{tabsndspeed}.

\begin{table}[t]
\newcommand{\tc}[1]{\textit{#1}}
\newcommand{\Z}{\hphantom{0}}
\caption[Rychlosti zvuku, hustoty, akustické impedance a absorpce]%
{\label{tabsndspeed}Rychlost šíření zvuku, hustota hmotnosti,
  akustická impedance a absorpce pro některá medicínsky důležitá prostředí.
  Hodnoty jsou pouze orientační a platí pro teplotu lidského organismu.
  Zpracováno podle~\cite{Dow-Ken-Joh:2006:PhysDiagImag:}.}
\begin{center}
\footnotesize
\begin{tabular}{lcccc}
\hline
\rule{0pt}{2.75ex}%
\tc{Prostředí}&\tc{Rychl. zvuku}&\tc{Hustota}&\tc{Akustická impedance}
&\tc{Absorpce} ($1\,\hz[M]$)\\
\rule[-1.75ex]{0pt}{0pt}%
              &$c\,\mathrm{[m\us s^{-1}]}$&$\rho\,\mathrm{[kg\us m^{-3}]}$
              &$Z\,\mathrm{[10^{6}\,kg\us m^{-2}\us s^{-1}]}$
  &$\alpha\,\mathrm{[dB\us cm^{-1}]}\,$\\\hline
Vzduch        & \Z330  & \Z\Z\Z\Z\hphantom{$.$}1.3 & \Z\Z\Z0.00043 & -- \\
Tuk           &  1470  & \Z970     &        1.42   & 0.6\Z  \\
Ricin. olej   &  1500  & \Z933     &        1.40   & --   \\
Voda          &  1492  &  1000     &        1.48   & -- \\
Měkká tkáň    &  1500  & \llap{$<{}$}1000 &   \llap{$\sim{}$}1.45 & 1.0\Z \\
Mozek         &  1530  &  1020     &        1.56   & 0.85 \\
Krev          &  1570  &  1020     &        1.60   & 0.18 \\
Ledviny       &  1561  &  1030     &        1.61   & -- \\
Játra         &  1549  &  1060     &        1.64   & 0.9\Z \\
Svaly         &  1568  &  1040     &        1.63   & 1.2 (3.3)$^{\ast}$ \\
Oční čočka    &  1620  &  1130     &        1.83   & 2.0\Z \\
Kost          &  4080  &  1700     &        6.12   & 20.0\Z\Z \\
\hline
\end{tabular}
\end{center}
{\footnotesize\mbox{}\enspace$^{\ast}$Podél (napříč) svalových vláken.}
\end{table}

\begin{figure}[t]
\begin{center}

\includegraphics[width=.72\linewidth]{longwav}
\end{center}
\caption[Ilustrace podélného vlnění]%
{\label{longwav}Místa zhuštění a zředění (a~jim odpovídající malé změny
  výchylek a rychlostí molekul ve směru šíření) se šíří zleva doprava (nebo
  obráceně) fázovou rychlostí $c$ jako podélná vlna. Ačkoli samy molekuly
  prostředí konají pouze malý oscilační pohyb okolo rovnovážné polohy, vlna
  sama přenáší energii a hybnost. Levý (pravý) panel reprezentuje vlnu
  s~vysokou (nízkou) amplitudou a intenzitou. Převzato
  z~\cite{Dow-Ken-Joh:2006:PhysDiagImag:}.}
\end{figure}

Prostředí, v~němž se vlna šíří, není bezztrátové; energie nesená vlnou je
disipována na teplo a intenzita, amplituda a další veličiny spojené s~energií
vlny (tlakové a hustotní změny, oscilační rychlost) klesají.
Obrázek~\ref{longwav} zobrazuje rovinnou vlnu, jejíž plochy konstantní
amplitudy (vlnoplochy) jsou roviny kolmé na směr šíření; v~případě sférické
vlny šířící se z~určitého bodu mají vlnoplochy sférický tvar a jejich povrch
se zvětšuje s~druhou mocninou vzdálenosti. Intenzita vlny proto v~důsledku
zákona zachování energie musí klesat jako převrácená hodnota kvadrátu
vzdálenosti, amplituda a akustický tlak musí klesat nepřímo úměrně vzdálenosti
(viz dále).

Libovolný průběh akustické výchylky lze lineárně rozložit pomocí
\emph{Fourierovy transformace} do harmonických složek popsaných funkcemi sinus
a kosinus; stačí proto zabývat se vlnou s~harmonickým průběhem šířící se ve
směru osy $x$
\begin{equation}
  u(x,t) = u_{\mathrm{max}}
    \cos 2\pi f\left(t-\frac{x}{c}\right)\,,                   \label{accu}
\end{equation}
kde $u_{\mathrm{max}}$ je amplituda výchylky z~rovnovážné polohy.  Okamžitá
\emph{akustická rychlost} kmitající částice prostředí podléhá zákonitosti
podobné~(\ref{accu}), ovšem s~amplitudou rychlosti $v_{\mathrm{max}}$, pro niž
derivací~(\ref{accu}) podle času dostaneme
\begin{gather}
  v_{\mathrm{max}} = 2\pi f u_{\mathrm{max}}\,.               \label{vampl}
\end{gather}
Typická velikost je řádově několik desítek milimetrů za sekundu.  Akustická
vlna je provázena, kromě výchylek částic prostředí~(\ref{accu}) a oscilací
rychlosti s~amplitudou~(\ref{vampl}), také oscilací tlaku kolem střední
hodnoty s~amplitudou~\cite{Brd-Sam-Sop:2005:MechKont:3}
\begin{equation}
  p_{\mathrm{max}} = \rho c v_{\mathrm{max}}\,,               \label{pampl}
\end{equation}
kde $\rho$ je hustota hmotnosti a $c$ rychlost zvuku v~prostředí.\footnote{Ve
  skutečnosti ve vlně osciluje také tlak s~amplitudou
  $\rho_{\mathrm{max}}\ll\rho$ okolo rovnovážné hodnoty $\rho$. Součin
  amplitud $\rho_{\mathrm{max}}v_{\mathrm{max}}$ je malá veličina druhého
  řádu, kterou zanedbáme vzhledem k~veličině prvního řádu $\rho
  v_{\mathrm{max}}$.}  Oscilační změny tlaku s~amplitudou~(\ref{pampl})
nazýváme \emph{akustickým tlakem}. Typická hodnota je několik
setin~$\mathrm{MPa}$.

Okamžitá intenzita zvuku je zvuková energie prošlá za jednotku času
jednotkovou plochou kolmou ke směru šíření ve fixním místě se souřadnicí $x$,
$I(t) = \rho c v^{2}(t)$.  Taková intenzita však není příliš užitečná; její
hodnota kolísá mezi nulou a maximální hodnotou odpovídající kladným i~záporným
extrémům rychlosti $\pm v_{\mathrm{max}}$. Místo ní se zavádí intenzita $I$
vystředovaná přes jednu periodu $T=1/f$; protože střední hodnota z~kvadrátu
sinu či kosinu je rovna $1/2$, je taková vystředovaná intenzita rovna
\begin{equation}
  I = \tfrac{1}{2} \rho c v_{\mathrm{max}}^{2}
    = \tfrac{1}{2} p_{\mathrm{max}} v_{\mathrm{max}}\,.      \label{aveint}
\end{equation}
Zavedeme-li \emph{efektivní hodnoty} akustické rychlosti a tlaku
\begin{equation}
  v_{\mathrm{ef}} \equiv \frac{v_{\mathrm{max}}}{\sqrt{2}}\,,\qquad
  p_{\mathrm{ef}} \equiv \frac{p_{\mathrm{max}}}{\sqrt{2}}\,,\label{efaccu}
\end{equation}
můžeme vztah pro střední intenzitu~(\ref{aveint}) napsat ve tvaru
\begin{equation}
    I = \rho c v_{\mathrm{ef}}^{2}
      = \frac{p_{\mathrm{ef}}^{2}}{\rho c}
      = p_{\mathrm{ef}} v_{\mathrm{ef}}\,.                  \label{accuint}
\end{equation}
Jednotkou intenzity je v~soustavě SI $\mathrm{W\us m^{-2}}$, avšak praktickou
jednotkou pro ultrazvukovou sonografii je $\mathrm{mW\us mm^{-2}}$. Typické
intenzity pro medicínský ultrazvuk jsou mezi $0.01\,\mathrm{mW\us mm^{-2}}$ a
$1\,\mathrm{mW\us mm^{-2}}$~\cite{Dow-Ken-Joh:2006:PhysDiagImag:}. Akustická
intenzita není na průřezu ultrazvukovým polem konstantní, nýbrž má maximum
uprostřed (SP${}={}$spatial peak) a průměrnou intenzitu označovanou jako
SA${}={}$spatial average intensity.

Pro porovnání dvou intenzit $I_{1}$, $I_{2}$ se obvykle používá logaritmické
škály v~\emph{decibelech} definované pomocí dekadických logaritmů jako
\begin{equation}
  L = 10\log\frac{I_{2}}{I_{1}}\,.                          \label{decibel}
\end{equation}
Označuje-li například $I_{2}$ intenzitu odražené vlny rovnou $0.0001$ násobku
intenzity dopadající vlny $I_{1}$, je $L=-40\,\mathrm{dB}$.
\emph{Polotloušťka} materiálu je tloušťka, na které intenzita klesne na
polovinu; a protože $\log0.5\approx0.3$, říkáme, že klesne
o~$3\,\mathrm{dB}$. Další užitečné hodnoty pro logaritmickou míru
v~$\mathrm{dB}$ jsou v~tabulce~\ref{tabdecibel}.

\begin{table}[t]
\caption[Tabulka často užívaných hodnot decibelů]%
{\label{tabdecibel}Tabulka často užívaných přibližných hodnot decibelů.
  Například změna
  o~$+3\,\mathrm{dB}$ znamená násobení faktorem${}\sim\!2$, změna
  o~$-3\,\mathrm{dB}$ znamená násobení faktorem ${}\sim\!1/2$. Jiné
  hodnoty než jsou
  uvedeny v~tabulce vhodně rozložíme, např. změna o $15\,\mathrm{dB}$
  rozložíme do změn o $10\,\mathrm{dB}$ a $5\,\mathrm{dB}$ a vynásobíme
  odpovídající faktory $10\times3=30$.}
\newcommand{\tc}[1]{\textit{#1}}
\newcommand{\Z}{\hphantom{0}}
\begin{center}
\small
\begin{tabular}{cc@{\hspace{4em}}cc}
\hline
\rule{0pt}{2.75ex}%
dB       & faktor & dB        & faktor  \\\hline
$\pm3$   &  $2$   & \Z$\pm7$  &  $\Z5$  \\
$\pm5$   &  $3$   & \Z$\pm9$  &  $\Z8$  \\
$\pm6$   &  $4$   &  $\pm10$  &  $10$   \\
\hline
\end{tabular}
\end{center}
\end{table}

Na \emph{útlumu ultrazvuku} se podílejí dva procesy: \emph{absorpce} v~tkáni,
jež tvoří hlavní mechanismus ($80\,\%$), a \emph{divergence ultrazvukového
  pole} způsobená odrazem a rozptylem. Frekvenční závislost absorpce je dána
viskozitou média, relaxačním časem a frekvencí ultrazvuku, a má typickou
závislost znázorněnou na obrázku~\ref{usabs} vlevo. Závislost intenzity na
dráze $s$ uražené v~homogenní tkáni je analogická exponenciálnímu zákonu pro
X-záření\INDEX{záření!X}
\begin{equation}
  I(s) = I_{0}\,\eto{-\beta s}\,,                            \label{usgabs}
\end{equation}
kde $I_{0}$ je počáteční intenzita ($s=0$), $\beta$ je \emph{absorpční
  koeficient}.  Místo absorpčního koeficientu $\beta$ se často udává
\emph{koeficient zeslabení} $\alpha$ definovaný jako
\begin{equation}
  \alpha
    =10\log\left.\frac{I_{0}}{I(s)}\right|_{s=1\,\mathrm{cm}}\,,\label{usatt}
\end{equation}
a udávaný v $\mathrm{dB\us cm^{-1}}$. Absorpční koeficient $\beta$ a
koeficient zeslabení $\alpha$ jsou frekvenčně závislé; pro některá běžná
medicínská prostředí při frekvenci $1\,\hz[M]$ jsou uvedeny
v~tabulce~\ref{tabsndspeed}.  Důležitým experimentálním faktem je, že mezi
koeficientem absorpce $\alpha$ a frekvencí je v~medicínsky relevantním rozsahu
ultrazvukových frekvencí přibližně lineární závislost; jinými slovy poměr
$\alpha/f$ je ve frekvenčním rozsahu $2.5$--$40\,\hz[M]$ přibližně konstantní.
Jednotlivé tkáně jsou proto charakterizovány hodnotou poměru $\alpha/f$
udávaného v~$\mathrm{dB\us cm^{-1}\us\hz[M]^{-1}}$, jejíž typická hodnota pro
měkké tkáně je zhruba $\sim\!1\mathrm{dB\us cm^{-1}\us\hz[M]^{-1}}$ (viz
obrázek~\ref{usabs} vpravo).

\begin{figure}[t]
\begin{center}
\includegraphics[width=.8\linewidth]{usabs}
\end{center}
\caption[Frekvenční závislost absorpce ultrazvuku v~tkáni]%
{\label{usabs} \emph{Vlevo:} Frek\-ve\-n\-č\-ní závislost absorpce ultrazvuku
  pro tkáň s~relaxační frekvencí $\omega$. \emph{Vpravo:} Závislost
  polotloušťky tkáně způsobující $50\,\%$ ztrátu intenzity na frekvenci pro
  typickou hodnotu $\alpha/f\sim1\mathrm{dB\us cm^{-1}\us\hz[M]^{-1}}$.
  \mbox{Převzato z~\cite{Dow-Ken-Joh:2006:PhysDiagImag:}}.}
\end{figure}

Je-li útlum ultrazvuku v~tkáni pro zobrazování nežádoucí jev, pak odraz
(reflexe) je naopak jev, jenž je podstatou ultrazvukového zobrazování.
V~akustice se pro popis reflexe používá \emph{akustické impedance $Z$},
veličiny definované jako
\begin{equation}
  Z = \frac{p_{\mathrm{ef}}}{v_{\mathrm{ef}}}
    = \rho c = \sqrt{K\rho}\,,                            \label{acuuimped}
\end{equation}
kde $K$ je modul pružnosti prostředí~\cite{Brd-Sam-Sop:2005:MechKont:3}.  Je
zkonstruována z~akustického tlaku a akustické rychlosti analogickým způsobem
jako elektrický odpor z~napětí a proudu v~Ohmově zákonu. Typické hodnoty pro
některé medicínsky relevantní tkáně jsou v~tabulce~\ref{tabsndspeed}.
Akustická impedance hraje důležitou roli při popisu chování ultrazvuku na
rozhraní dvou různých tkání.

Rozlišujeme \emph{zrcadlový odraz} na hladkém rozhraní\footnote{Hladkým
  rozhraním rozumíme rozhraní, jehož nepravidelnosti jsou velmi malé ve
  srovnání s~vlnovou délkou $\lambda$ ultrazvukového svazku; např. pro
  frekvenci $2\,\hz[M]$ je v~měkké tkáni $\lambda=0.75\,\mathrm{mm}$.} mezi
dvěma tkáněmi s~různou akustickou impedancí. Pro kolmý dopad je zrcadlový
odraz znázorněn na obrázku~\ref{usreflex} vlevo. Vztah pro \emph{reflexivitu}
a \emph{transmisivitu} lze pomocí akustické impedance psát ve tvaru (viz
obrázek~\ref{usreflex} pro význam veličin)
\begin{gather}
  R\equiv \frac{I_{\mathrm{r}}}{I_{\mathrm{i}}}
    = \left(\frac{Z_{1}-Z_{2}}{Z_{1}+Z_{2}}\right)^{2}\,,    \label{usiref}\\
  T\equiv \frac{I_{\mathrm{t}}}{I_{\mathrm{i}}}
    = \frac{4Z_{1}Z_{2}}{(Z_{1}+Z_{2})^{2}}\,.             \label{usitrans}
\end{gather}
Snadno se ověří, že $R+T=1$, jak vyžaduje splnění zákona zachování energie.
Při průchodu rozhraním vzduch-tuk s~podstatně odlišnými akustickými
impedancemi se podle údajů v~tabulce~\ref{tabsndspeed} a vztahů (\ref{usiref})
a (\ref{usitrans}) odrazí $R = 99.9\,\%$ intenzity (jedná se o~\emph{echogenní
  strukturu}), zatímco při průchodu rozhraním játra-ledviny s~velmi podobnými
akustickými impedancemi se odrazí pouze $R = 0.0085\,\%$ intenzity a $T =
99.9915\,\%$ intenzity prochází.

\begin{figure}[t]
\begin{center}

\includegraphics[width=\linewidth]{usreflex}
\end{center}
\caption[Odraz ultrazvukového vlnění]%
{\label{usreflex} \emph{Vlevo:} Odraz ultrazvukového svazku dopadajícího kolmo
  na hladké rozhraní mezi dvěma tkáněmi o~akustických impedancích $Z_{1}$ a
  $Z_{2}$. Intenzita dopadajícího svazku $I_{\mathrm{i}}$ se rozdělí mezi
  prošlý svazek $I_{\mathrm{t}}$ a odražený svazek $I_{\mathrm{r}}$. Rozdělení
  intenzit závisí na akustických impedancích tkání $Z_{1}$ a $Z_{2}$ podle
  vztahů~(\ref{usiref}) a~(\ref{usitrans}). \emph{Vpravo:} Při dopadu pod
  nenulovým úhlem $\theta_{\mathrm{i}}$ dochází ke změně směru šíření prošlého
  svazku. Pro úhly lomu $\theta_{\mathrm{t}}$ a odrazu $\theta_{\mathrm{r}}$
  platí analogické zákony jako v~optice, zákon odrazu~(\ref{usrelaw}) a zákon
  lomu~(\ref{usisnell}).  Převzato z~\cite{Dow-Ken-Joh:2006:PhysDiagImag:}.}
\end{figure}

Při průchodu rozhraním s~nepravidelnostmi nebo strukturou o~velikosti
srovnatelné s~vlnovou délkou ultrazvuku, nebo při průchodu granulózním
prostředím (krevní částice), dochází k~\emph{difrakčnímu jevu}~--
\emph{rozptylu}.  Rozptýlené ultrazvukové vlnění formuje kužel, jehož
vrcholový úhel závisí nepřímo úměrně na vlnové délce $\lambda$ a přímo úměrně
na charakteristické velikostní škále nepravidelností rozhraní.  Rozptyl se
také zvětšuje s~rostoucím úhlem dopadu a jeho intenzita je úměrná $f^{4}$.
Rozptyl nemusí být nutně jen parazitním jevem; charakteristické interferenční
obrazce nazývané \emph{speckle} totiž svým tvarem odpovídají určité morfologii
rozhraní a mohou tak detekovat patologii projevující se právě změnou
morfologie rozhraní~-- příkladem je diagnostika\INDEX{diagnostika!onemocnění
  jater} onemocnění jater.  Difrakce se také projevuje obrazovými artefakty
velikosti řádově vlnové délky použitého ultrazvuku při průchodu svazku kolem
náhlých akustických stínění jako například~kosti.

Kromě kolmého odrazu se při průchodu rozhraním pod nenulovým úhlem dopadu
$\theta_{\mathrm{i}}$ uplatňují zákony analogické zákonům geometrické optiky
(viz obrázek~\ref{usreflex} vpravo), konkrétně známý \emph{zákon odrazu}
\begin{equation}
  \theta_{\mathrm{r}} = \theta_{\mathrm{i}}                 \label{usrelaw}
\end{equation}
a \emph{Snellův zákon lomu}
\begin{equation}
  \frac{\sin\theta_{\mathrm{i}}}{v_{\mathrm{i}}}
    = \frac{\sin\theta_{\mathrm{t}}}{v_{\mathrm{t}}}\,,    \label{usisnell}
\end{equation}
kde $v_{\mathrm{i}}$ ($v_{\mathrm{t}}$) je fázová rychlost ultrazvuku
v~prostředí dopadajícího (prošlého) svazku.

\subsection{Ultrazvuková sonda}\label{sontransd}%//////////////////////////

Ultrazvuková sonda (též ultrazvukový transducer) přeměňuje energii
elektrických pulsů přiváděných z~generátoru do sondy na energii ultrazvukových
vln vysílaných do vyšetřovaného objektu, a~opačně, dokáže přeměňovat energii
ultrazvukových odrazů přicházejících z~nitra vyšetřovaného objektu na dále
zpracovatelný elektrický signál. Je obvykle založena na \emph{piezoelektrickém
  jevu}, jenž umožňuje uvedený \uv{oboustranný} převod. Základní konstrukční
schéma piezoelektrické ultrazvukové sondy je na obrázku~\ref{ustransd}.

\begin{figure}[t]
\begin{minipage}{.39\linewidth}
\centering

\includegraphics[width=0.8\linewidth]{ustransda}\par

\includegraphics[width=\linewidth]{ustransdb}
\end{minipage}\hfill%
\begin{minipage}{.57\linewidth}
\caption[Základní konstrukce ultrazvukové sondy]%
{\label{ustransd}Schéma konstrukce ul\-tra\-zvu\-ko\-vé sondy.
  \emph{Nahoře:} Základem je piezoelektrický krystal (obvykle
  $\mathrm{Pb\,Zr\,Ti\,O_{4}}$) o~tloušťce optimálně $\lambda/2$, pokrytý
  přizpůsobovací vrstvou (Matching plastic layer). Pro optimální přenos
  ultrazvuku z~piezokrystalu by akustická impedance $Z_{\mathrm{M}}$ této
  vrstvy měla být geometrickým průměrem
  $Z_{\mathrm{M}}=\sqrt{Z_{\mathrm{T}}Z_{\mathrm{L}}}$ akustické impedance
  piezokrystalu $Z_{\mathrm{T}}$ a tkáně $Z_{\mathrm{L}}$, a~její tloušťka by
  měla být rovna čtvrtině vlnové délky. Za piezokrystalem se nachází tlumicí
  blok (Backing block) z~jemných wolframových částeček suspendovaných
  v~epoxidové pryskyřici (wolframové částečky rozptylují ultrazvuk, jenž je
  pak v~pryskyřici pohlcován). \emph{Dole:} Mezi sondu a pokožku či jiný
  povrch vyšetřovaného pacienta se nanáší tenká vrstva kontaktního gelu.
  Převzato z~\cite{Dow-Ken-Joh:2006:PhysDiagImag:}.}
\end{minipage}
\end{figure}

\begin{figure}[t]
\begin{center}

\includegraphics[width=\linewidth]{uspulse}
\end{center}
\caption[Geometrie ultrazvukových pulsů]%
{\label{uspulse}Geometrie ultrazvukových pulsů. \emph{Vlevo:} Ultrazvukový
  puls obsahující tři cykly o~vlnové délce $\lambda$ s~vhodně tvarovanou
  obálkou. Zkratka SPL značí \emph{Spatial Pulse Length} (prostorovou délku
  pulsu), typická hodnota je $0.45\,\mathrm{mm}$ až $2.25\,\mathrm{mm}$.
  \emph{Vpravo:} Podobně PD značí \emph{Pulse Duration} neboli trvání pulsu,
  které bývá v~rozmezí $0.5\,\micro\mathrm{s}$ až $3\,\micro\mathrm{s}$, PRP
  značí \emph{Pulse Repetition Period} neboli opakovací periodu pulsů, která
  bývá v~řádu $\mathrm{ms}$, případně $\mathrm{PRF}=1/\mathrm{PRP}$
  \emph{Pulse Repetition Frequency} značí opakovací frekvenci pulsů.  Převzato
  z~\cite{Dow-Ken-Joh:2006:PhysDiagImag:}.}
\end{figure}

Ultrazvukové pulsy vytvářené sondou jsou velmi krátké~-- obvykle obsahují jen
dvě až tři vlny. Typický tvar pulsu je na obrázku~\ref{uspulse}.  Z~teorie
Fourierovy transformace plyne, že čím více zkrátíme trvání pulsu (PD), tím je
spektrum širší. Záleží také na obálce (tvaru) pulsu~-- tvar by se měl alespoň
přibližovat Gaussově křivce, aby byly potlačeny oscilační laloky vzdálené od
základní frekvence. Opakovací frekvence pulsů (PRP) je řádově tisíc krát větší
než trvání pulsu (PD), protože po vyslání každého pulsu se sonda přepne do
\uv{naslouchacího} režimu, v~němž registruje ultrazvuková echa přicházející
z~nitra vyšetřovaného pacienta, než opět vyšle krátký puls. Tento cyklus se
neustále opakuje. Vzhledem k~proporcím mezi PD a PRP je sonda zhruba
$99.9\,\%$ času v~naslouchacím režimu.

Opakovací frekvence PRF jednoduše souvisí s~maximální hloubkou obrazu $D$: čas
potřebný na proniknutí ultrazvuku do hloubky $D$ a na návrat zpět do sondy je
$2D/c$. Aby sonda byla ještě schopna detekovat příslušnost echa k~vyslanému
pulsu, musí být opakovací perioda PRP delší než $2D/c$. Z~této úvahy dostáváme
\begin{equation}
  \mathrm{PRF} < \frac{c}{2D}\,.                            \label{usimdep}
\end{equation}
Pro hloubku penetrace $150\,\mathrm{mm}$ a
$c\approx1500\,\mathrm{m\us s^{-1}}$ máme $\mathrm{PRF}\approx
5\,\hz[k]$. Typické hodnoty pro ultrazvukovou sonografii jsou
$\mathrm{PRF}\approx 2$--$10\,\hz[k]$. Opakovací frekvence musí být adekvátní
hloubce obrazu, příliš krátká PRP může způsobit detekci echa z~předchozích
pulsů a působit falešná zobrazení.

Frekvenční spektrum ultrazvukových pulsů má tvar Gaussovy křivky; faktor
kvality piezokrystalu $Q$ je dán poměrem rezonanční frekvence
$f_{\mathrm{R}}$, na níž má tato křivka maximum, a~šířky křivky $\Delta f$
(rozdílu frekvencí, pro něž spektrální křivka poklesne na
$1/\sqrt{2}\approx0.707$ maxima $f_{\mathrm{R}}$, viz obrázek~\ref{uslobes}
vlevo):
\begin{equation}
  Q = \frac{f_{\mathrm{R}}}{\Delta f}\,.                    \label{usqualf}
\end{equation}
Sondy s~nižším $Q$ (tj. širším spektrem a krátkými pulsy) se používají pro
ultrazvukové zobrazování, kde potřebujeme pro dobré rozlišení co nejkratší SPL
pro co nejvyšší \emph{axiální rozlišovací schopnost} (ta se obvykle definuje
jako $\mathrm{SPL}/2$, je tedy typicky $0.2\,\mathrm{mm}$ až
$1.2\,\mathrm{mm}$) a nezáleží příliš na přesné frekvenci. Naopak sondy
s~vyšším $Q$ se používají v~dopplerovské sonografii (viz níže), kde záleží na
přesné frekvenci a délka trvání pulsu může být velká (nebo kontinuální).

\begin{figure}[p]
\begin{center}

\includegraphics[width=.86\linewidth]{uslobes}
\end{center}
\caption[Faktor kvality a potlačení postranních laloků]%
{\label{uslobes}\emph{Vlevo:} Rozdíl $\Delta f$ frekvencí, pro něž spektrální
  křivka poklesne na $1/\sqrt{2}\,f_{\mathrm{R}}\approx0.707f_{\mathrm{R}}$.
  Spektrum sondy s~vysokým $Q$ (malým $\Delta f = B$) je vyznačeno černě,
  spektrum sondy s~nízkým $Q$ (vysokým $\Delta f = A$) je vyznačeno šedě.
  \emph{Vpravo:} Potlačení postranních laloků pomocí apodizace.  Převzato
  z~\cite{Dow-Ken-Joh:2006:PhysDiagImag:}.}
\end{figure}
\begin{figure}[p]
\begin{center}

\includegraphics[width=.86\linewidth]{uszones}
\end{center}
\caption[Ilustrace Fresnelovy a Fraunhoferovy zóny]%
{\label{uszones}Ilustrace změny Fresnelovy a Fraunhoferovy zóny v~závislosti
  na rozměru ultrazvuk produkujícího elementu sondy a na frekvenci. Převzato
  z~\cite{Dow-Ken-Joh:2006:PhysDiagImag:}.}
\end{figure}
\begin{figure}[p]
\begin{center}

\includegraphics[width=\linewidth]{usfocus}
\end{center}
\caption[Vliv fokusace sondy na tvar ultrazvukového pole]%
{\label{usfocus}Vliv fokusace sondy na tvar ultrazvukového pole.  Zleva
  doprava: plochá sonda, slabě fokusovaná ($\ell_{\mathrm{near}}/R<2$) a silně
  fokusovaná ($\ell_{\mathrm{near}}/R\approx3$--$4$). Veličina $\alpha$ na
  obrázku odpovídá poloměru sondy $a$. Převzato
  z~\cite{Dow-Ken-Joh:2006:PhysDiagImag:}.}
\end{figure}

Tvar akustické apertury (profil ultrazvukového pole) je vhodné upravit tak,
aby víceméně kopíroval Gaussovu křivku. Při obdélníkovém profilu bychom totiž
dostali, jak už to u~zdrojů vlnění bývá, kromě hlavního laloku navíc
\emph{postranní laloky}, které by způsobovaly degradaci kontrastu a rozlišení.
Přizpůsobení akustické apertury se nazývá \emph{apodizace} a její výsledek je
na obrázku~\ref{uslobes} vpravo. Apodizace se dosahuje tvarovanými akustickými
čočkami a dynamickou fokusací (viz níže a obrázek~\ref{usfocus}).

Na závěr si všimneme chování ultrazvukového pole podél směru šíření, jež má
vliv na \emph{stranovou (laterální) rozlišovací schopnost}. Pole lze
rozdělit do dvou oblastí: na \emph{blízkou (Fresnelovu) zónu} a na
\emph{vzdálenou (Fraunhoferovu) zónu} (viz obrázek~\ref{uszones}). Zatímco
ve Fresnelově zóně o~délce
\begin{equation}
  \ell_{\mathrm{near}} = \frac{a^{2}}{\lambda}\,,         \label{usfresnel}
\end{equation}
kde $a$ je poloměr sondy, zachovává pole šířku sondy, ve vzdálené zóně se
stává divergentním s~úhlem divergence
\begin{equation}
  \theta_{\mathrm{d}}
    = \arcsin\left(\frac{0.61\lambda}{a}\right)\,.         \label{usdivang}
\end{equation}
Pro lepší laterální rozlišení ve velkých hloubkách potřebujeme širší sondu a
vysokou frekvenci, aby blízká zóna byla co nejdelší. V~praxi je však nutné
volit kompromisní hodnoty a přihlížet k~dalším aspektům šíření ultrazvuku
v~tkáni, zejména k~faktu, že s~rostoucí frekvencí roste útlum. Například
v~oftalmologii stačí jen velmi mělká hloubka a sonda má rovněž malý rozměr,
takže se používá ve srovnání s~ostatními aplikacemi ultrazvukového zobrazování
vyšších frekvencí.

Laterální rozlišení lze ovlivnit vhodným tvarováním (fokusací) sondy, jak je
vidět z~obrázku~\ref{usfocus}.

\section{Nukleární magnetická rezonance (NMR/MRI)}\label{nmr}%%%%%%%%%%%%%%

Nukleární magnetická rezonance (NMR) je fyzikálně-chemická metoda založená na
analýze magnetických momentů atomových jader. Tato původně analytická metoda
objevená \designation{F. Blochem} a \designation{E. Purcellem} v~roce~1946
byla pod názvem MRI (Magnetic Resonance Imaging) rozvinuta jako zobrazovací
technika používaná k~tomografickému zobrazování vnitřních orgánů lidského
těla. Obrazy vznikají na základě odezvy dočasných změn magnetických poměrů
v~tkáních. Hned zkraje poznamenejme, že v~metodě MRI nefiguruje ionizující
záření\INDEX{záření!ionizující} jako v~klasické či digitální rentgenologii
nebo v~počítačové tomografii; využívá se radiofrekvenčního elektromagnetického
záření\INDEX{záření!radiofrekvenční}. Při zpracování této sekce byly převážně
využity zdroje~\cite{Hor:LIVE:BasicsMRI:,%
  Ulm:LIVE:AplIonZarJRM:,%
  Dow-Ken-Joh:2006:PhysDiagImag:,%
  Hau:1990:CZEJPA:,%
  Cer:2004:POKMF:NMRImag}.

\subsection{Stručný nástin principu MRI}\label{mribrief}%//////////////////

Atomové jádro se skládá z~protonů a neutronů, jež souhrnně nazýváme nukleony.
Každý nukleon má vlastní moment hybnosti~-- spin. V~jádrech s~lichým počtem
nukleonů se spiny jednotlivých nukleonů nekompenzují a určitý spin nese celé
jádro. Pro lékařské účely je důležité především jádro vodíku, tedy
proton.\footnote{Z~toho důvodu se MRI užíval i~název \emph{vodíkové
    topografické zobrazení}~\cite{Ulm:LIVE:AplIonZarJRM:}.} Spin je kvantovaná
veličina, může být vzhledem k~určitému směru kvantování orientována buď
\uv{paralelně} nebo \uv{antiparalelně}.

Mějme soubor jader vodíku. Pokud nejsou vložená do vnějšího magnetického pole,
bude polovina protonů orientovaná nahoru a polovina dolů a jejich energie bude
stejná. Přiložíme-li vnější magnetické pole, bude energie \uv{elementárních
  magnetů}, jaké protony představují, různá pro protony orientované ve směru a
proti směru pole. Dopadne-li nyní na tento vzorek vysokofrekvenční
elektromagnetické záření\INDEX{záření!elektromagnetické} obsahující fotony
s~energií rovné rozdílu potenciální energie protonů s~antiparalelní a
paralelní orientací, dojde k~rezonančnímu jevu: foton může být absorbován a
jeho energie excituje proton v~nižším energetickém stavu do vyššího
energetického stavu. Pokud záření dále nepůsobí, dojde k~relaxaci, proton se
opět \uv{přepne} do nižšího energetického stavu a přitom se uvolní foton. Při
rezonanční frekvenci dochází ke zvýšené absorpci elektromagnetického vlnění
v~látce, která je registrována prostřednictvím fotonů uvolňovaných
při~relaxaci.

\begin{figure}[t]
\begin{minipage}{.19\linewidth}
\centering
\includegraphics[width=\linewidth]{ive1}
\end{minipage}\hfill%
\begin{minipage}{.19\linewidth}
\centering
\includegraphics[width=\linewidth]{ive2}
\end{minipage}\hfill%
\begin{minipage}{.19\linewidth}
\centering
\includegraphics[width=\linewidth]{joe3}
\end{minipage}\hfill%
\begin{minipage}{.19\linewidth}
\centering
\includegraphics[width=\linewidth]{joe4}
\end{minipage}\hfill%
\begin{minipage}{.19\linewidth}
\centering
\includegraphics[width=\linewidth]{joe5}
\end{minipage}
\caption[Lokalizace MRI řezu pomocí gradientního pole]%
{\label{mriloc}Silné magnetické pole používané k~MRI zobrazování má
  \emph{lineární gradient}, tj. mění se rovnoměrně v~určitém směru (na obrázku
  ve směru rovnoběžném se hřbetem vazby).  K~rezonanci dochází pro frekvenci
  odpovídající vybranému typu jader jen při určité intenzitě magnetického
  pole, jež leží v~tenké rovinné vrstvě (řezu) o~tloušťce Thk kolmé ke směru
  gradientu. Pomocí gradientu magnetického pole se tak dá vybírat řez
  pacientem, aniž by bylo nutné s~pacientem fyzicky manipulovat.  Zobrazovaná
  vrstva je složena z~objemových \emph{voxelů} (trojrozměrný analogon pixelu,
  obvyklý objem voxelu je${}\sim\!3\,\mathrm{mm^{3}}$). Zcela vpravo je
  dvourozměrné zobrazení vybraného řezu adresovaného odpovídajícím gradientem
  magnetického pole.  Upraveno podle~\cite{Hor:LIVE:BasicsMRI:}.}
\end{figure}

\begin{figure}[t]
\begin{minipage}{.16\linewidth}
\centering

\includegraphics[width=\linewidth]{intro_3i}
\end{minipage}\hfill%
\begin{minipage}{.16\linewidth}
\centering

\includegraphics[width=\linewidth]{intro_3j}
\end{minipage}\hfill%
\begin{minipage}{.16\linewidth}
\centering

\includegraphics[width=\linewidth]{intro_3l}
\end{minipage}\hfill%
\begin{minipage}{.16\linewidth}
\centering

\includegraphics[width=\linewidth]{intro_3m}
\end{minipage}\hfill%
\begin{minipage}{.16\linewidth}
\centering

\includegraphics[width=\linewidth]{intro_3n}
\end{minipage}\hfill%
\begin{minipage}{.16\linewidth}
\centering

\includegraphics[width=\linewidth]{intro_3o}
\end{minipage}
\caption[Ilustrace obsahu vodíkových jader v~tkáni]%
{\label{mricell}Lidské tělo obsahuje mnoho vody a lipidů. Voda i~lipidy
  obsahují mnoho vodíkových atomů (v $1\,\mathrm{cm^{3}}$ vody je řádově
  $10^{22}$ vodíkových atomů), takže lidské tělo obsahuje asi $63\,\%$ vodíku.
  Protože vodíkové atomy jsou schopné dávat NMR signál, používá se pro
  zobrazování právě jich.  Každý voxel obsahuje jednu nebo více tkání se
  specifickým zastoupením vodíku (zcela vlevo). Zoom voxelu odhalí buňky
  (druhý panel zleva).  Uvnitř každé buňky jsou vodní molekuly (třetí a čtvrtý
  panel zleva). Každá molekula vody obsahuje jeden kyslíkový a dva vodíkové
  atomy (druhý panel zprava).  Jádro vodíkového atomu je tvořeno protonem~--
  miniaturním \uv{magnetem}, jenž se využívá v~metodě NMR, potažmo
  k~medicínskému zobrazování. Právě různý obsah vodíkových atomů a různý
  chemismus obklopujícího prostředí v~různých tkáních umožňuje jejich
  rozlišení.  Převzato z~\cite{Hor:LIVE:BasicsMRI:}.}
\end{figure}

V~praktické aplikaci se vzorek umístí do silného magnetického pole, které
způsobí částečnou polarizaci protonů. Rezonanční frekvence závisí na intenzitě
magnetického pole. Jeli vzorek umístěn do nehomogenního pole, tedy pole, které
má určitý gradient, k~rezonanci dochází v~různých místech při různých
frekvencích, čímž je možno hledanou strukturu lokalizovat~(viz
obrázek~\ref{mriloc}). Prakticky všechny tkáně obsahují značné množství vody,
tedy vodíkových jader (viz obrázek~\ref{mricell}). Rozložení molekul vody
v~příslušné tkáni tedy může být zobrazeno popsaným způsobem. Pro použitelnost
metody je podstatné, že pomocí gradientu magnetického pole je možno přesně
lokalizovat polohu zdroje relaxačního signálu. V~praktickém provedení je
samozřejmě tento signál dále zesilován a zpravidla analyzován pomocí
Fourierovy transformace. Vhodným měněním gradientu přiloženého magnetického
pole je možné získávat tomografický obraz.

\subsection{Fyzika jaderného spinu}\label{nuclespin}%//////////////////////

Jak již bylo výše řečeno, každý nukleon (proton nebo neutron), z~nichž se
skládají atomová jádra, má \emph{spin}.\footnote{Spin je kvantovaná veličina a
  u~fermionů, mezi něž proton a neutron patří, nabývá poločíselného násobku
  redukované Planckovy konstanty $\hbar$. Proton i~neutron mají v~jednotkách
  $\hbar$ spin~$1/2$.}  Se spinem je svázán i~vlastní (spinový)
\emph{magnetický moment}. Spin a magnetický moment mají však pouze atomová
jádra s~lichým nukleonovým číslem $A$, neboť spiny a magnetické momenty
spárovaných protonů a neutronů se vzájemně ruší.  Magnetický moment jádra
vytváří nespárovaný nukleon. Magnetickou rezonanci lze tedy pozorovat jen
u~jader s~lichým nukleonovým číslem. Pro potřeby medicínského zobrazování je
ještě zapotřebí tuto podmínku doplnit o~podmínku dostatečného zastoupení
v~tkáních. Zdaleka největší význam má vodík ${}^{1}\mathrm{H}$, jehož průměrné
zastoupení v~tkáních je kolem $63\,\%$.  Další atomy, jejichž jádra přicházejí
pro medicínské zobrazování v~úvahu, jsou například izotop uhlíku
${}^{13}\mathrm{C}$, fosfor ${}^{31}\mathrm{P}$, izotop dusíku
${}^{15}\mathrm{N}$, fluor ${}^{19}\mathrm{F}$ a sodík ${}^{23}\mathrm{Na}$
(viz tabulka~\ref{tablarm}).

\begin{table}[p]
  \caption[Parametry izotopů pro medicínské zobrazování]%
  {\label{tablarm}Parametry izotopů pro medicínské zobrazování.
    Podle~\cite{Hor:LIVE:BasicsMRI:,Dow-Ken-Joh:2006:PhysDiagImag:}.}
\begin{center}
\small
\newcommand{\pad}{\hphantom{0}}
\begin{tabular}{lccc}
\hline
Izotop & Gyromagnetický & Abundance & Relativní inten- \\
       & poměr $\gamma$ [$\hz[M]\us\mathrm{T^{-1}}$] & [\%] &zita na proton\\\hline
${}^{1}\mathrm{H}$ (proton) & $42.58$ & \hphantom{0}99.98\hphantom{0}&1.0000 \\
${}^{19}\mathrm{F}$         & $40.05$ & 100\hphantom{.000} & 0.8300 \\
${}^{23}\mathrm{Na}$        & $11.26$ & 100\hphantom{.000} & 0.0930 \\
${}^{31}\mathrm{P}$         & $17.24$ & 100\hphantom{.000} & 0.0660 \\
${}^{17}\mathrm{O}$      & $\pad5.77$ & \hphantom{00}0.037  & 0.0290 \\
${}^{13}\mathrm{C}$         & $10.71$ & \hphantom{00}1.11\hphantom{0}&0.0160 \\
${}^{15}\mathrm{N}$      & $\pad4.30$ & \hphantom{00}0.37\hphantom{0}&0.0010 \\
\hline
\end{tabular}
\end{center}
\end{table}

\begin{figure}[p]
\begin{center}

\includegraphics[width=0.9\linewidth]{nucspin}
\end{center}
\caption[Chování jaderných spinů ve vnějším magnetickém poli]%
{\label{mrispin}Názorná (i~když z~hlediska kvantové mechaniky nepřesná)
  představa chování jaderných spinů ve vnějším magnetickém poli. \emph{Vlevo:}
  Spiny jsou bez vnějšího magnetického pole orientovány chaoticky a jejich
  celkový magnetický moment je nulový.  \emph{Uprostřed:} Zapneme-li vnější
  magnetické pole, např. vpuštěním elektrického proudu do cívky obklopující
  vzorek tkáně, budou se jaderné spiny snažit \uv{srovnávat} souhlasně
  s~magnetickým polem a začnou vytvářet makroskopický magnetický moment~--
  tkáň je \uv{zmagnetovaná}. \emph{Vpravo:} Aplikací elektromagnetického
  záření o~Larmorově frekvenci~(\ref{larmor}) dochází k~rezonanci a jaderné
  spiny s~nižší energií jsou excitovány do vyšší energetické hladiny, což
  ovlivňuje hodnotu celkové makroskopické magnetizace.  Převzato
  z~\cite{Ulm:LIVE:AplIonZarJRM:}.}
\end{figure}

\begin{figure}[p]
\begin{minipage}{.32\linewidth}
\centering

\includegraphics[width=\linewidth]{t1}
\end{minipage}\hfill%
\begin{minipage}{.32\linewidth}
\centering

\includegraphics[width=\linewidth]{mzt}
\end{minipage}\hfill%
\begin{minipage}{.32\linewidth}
\centering

\includegraphics[width=\linewidth]{mxyt}
\end{minipage}
\caption[Ilustrace ke spinové relaxaci]%
{\label{longrelax}Ke spin-mřížkové a spin-spinové relaxaci. \emph{Vlevo:}
  V~rovnováze leží vektor magnetizace $M_{0}$ podél vektoru vnějšího
  magnetického pole $B_{0}$. \emph{Uprostřed:} $T_{1}$ je čas, za který se
  podélná magnetizace vrátí z~nulové hodnoty na $1-\eto{-1}\approx 63\,\%$
  rovnovážné hodnoty $M_{0}$. \emph{Vpravo:} $T_{2}$ je čas, za který příčná
  magnetizace klesne z~vychýlené hodnoty $M_{\perp0}$ na
  $\eto{-1}\approx 37\,\%$ této hodnoty.  Převzato
  z~\cite{Hor:LIVE:BasicsMRI:}.}
\end{figure}

Bez vnějšího magnetického pole bude polovina protonů orientována jedním směrem
a polovina opačným, takže v~makroskopickém měřítku se jejich magnetické
momenty vyruší a makroskopický magnetický moment tkáně bude nulový.
Zapneme-li vnější magnetické pole o~indukci řádově několik Tesla,\footnote{To
  je velmi silné magnetické pole, řádově $10^{4}\times$ silnější než přirozené
  geomagnetické pole Země. Pro vytvoření takto silných magnetických polí se
  obvykle používají supravodivé cívky využívající chlazení kapalným héliem.}
každý spin se srovná do směru pole v~jedné ze dvou možných orientací,
paralelní s~nižší energií nebo antiparalelní s~vyšší energií.  Tkáň rychle
přejde do rovnovážného stavu, v~němž počet jader $N^{+}$ orientovaných
paralelně (s~nižší energií) bude mírně převyšovat počet jader $N^{-}$
orientovaných antiparalelně.\footnote{Statistická fyzika umožňuje spočítat
  jejich poměr, $N^{+}/N^{-} = \eto{\Delta E/kT}>1$, kde
  $\Delta E = h\gamma B_{0}$ je rozdíl energie antiparalelně a paralelně
  orientovaných spinů, $k$ je Boltzmannova konstanta a $T$ je absolutní
  teplota. Nadbytek paralelně orientovaných spinů je malý; rozdíl
  $N^{+}-N^{-}$ je při běžné teplotě lidského těla a magnetické indukci
  $1\,\mathrm{T}$ veličina řádu $10^{16}$ spinů na $1\,\mathrm{cm^{3}}$ látky,
  dostačuje však k~vytvoření měřitelné makroskopické magnetizace.} Součet
jaderných spinů v~určité oblasti se stejnou hodnotou vnějšího magnetického
pole $B_{0}$ bude nenulový a bude vytvářet makroskopickou \emph{magnetizaci}
tkáně $M_{0}$.  Velikost této magnetizace bude úměrná rozdílu $N^{+}-N^{-}$.
Tento proces \uv{zmagnetování} tkáně je názorně ilustrován na
obrázku~\ref{mrispin} vlevo a uprostřed. Proces nastavení tepelné rovnováhy
mezi spinovým systémem a zcela obecným okolím označovaným jako \emph{mřížka}
(i~v~případě kapalin) se nazývá \emph{spin-mřížková} nebo také \emph{podélná}
relaxace.

Nyní pomocí další cívky napájené vysokofrekvenčním elektrickým signálem
vyšleme do vzorku tkáně elektromagnetické
záření\INDEX{záření!elektromagnetické} obsahující fotony s~takovou frekvencí
$f$, aby se jejich energie $E$ podle Planckova vztahu $E=hf$ (kde $h$ je
Planckova konstanta) rovnala rozdílu mezi energií antiparalelně a paralelně
orientovaného spinu. Tehdy dochází k~jevu \emph{jaderné magnetické rezonance}:
foton může být absorbován a jeho energie excituje proton z~nižšího
(paralelního) do vyššího (antiparalelního) energetického stavu. Frekvence $f$
splňující rezonanční podmínku se nazývá \emph{Larmorovou frekvencí} a
s~indukcí vnějšího magnetického pole $B_{0}$ souvisí vztahem
\begin{equation}
  f = \gamma B_{0}\,,                                        \label{larmor}
\end{equation}
kde $\gamma=42.58\,\hz[M]\us\mathrm{T^{-1}}$ se nazývá \emph{gyromagnetický
  poměr} protonu. Další prvky, resp. jejich izotopy použitelné pro potřeby
medicínského zobrazování, shrnuje tabulka~\ref{tablarm}.

Je zřejmé, že popsaná excitace jaderných spinů v~režimu rezonance ovlivní
rovnovážnou magnetizaci tkáně $M_{0}$. V~rovnováze leží vektor magnetizace
$M_{0}$ podél vektoru vnějšího magnetického pole $B_{0}$ a je proto nazýván
rovnovážnou magnetizací (obrázek~\ref{longrelax} vlevo). Položíme-li osu $z$
kartézského souřadnicového systému do tohoto směru, pak v~rovnováze je podélná
komponenta magnetizace $M_{z}=M_{0}$ a příčná složka $M_{\perp}$ sestávající
z~komponent do os $x$, $y$ kolmých k~ose $z$ je rovna nule, $M_{\perp} = M_{x}
= M_{y} = 0$.  Dodáváme-li larmorovské elektromagnetické
záření,\INDEX{záření!elektromagnetické} jež excituje jaderné spiny, ve formě
vhodně dlouhého a tvarovaného pulsu\footnote{Nazývaného též RF puls,
  radiofrekvenční puls, jehož základní frekvence je např. pro magnetické pole
  $B_{0}=1\,\mathrm{T}$ rovna $42.58\,\hz[M]$, pro magnetické pole
  $B_{0}=1.5\,\mathrm{T}$ rovna $63.87\,\hz[M]$, podle vztahu~(\ref{larmor}).}
můžeme dosáhnout anulování podélné magnetizace, tj.  $M_{z}=0$. Po odeznění RF
pulsu se ovšem jaderné spiny vracejí zpět do původní rovnováhy~-- říkáme, že
\emph{relaxují}~-- a časový vývoj podélné magnetizace se chová podle rovnice
\begin{equation}
  M_{z}(t) = M_{0}\left(1-\eto{-\frac{t}{T_{1}}}\right)\,, \label{spinlat}
\end{equation}
kde $T_{1}$ se nazývá \emph{spin-mřížkový relaxační čas}. Jeho hodnota
charakterizuje rychlost návratu magnetizace zpět do rovnovážné hodnoty. Tento
jev se nazývá \emph{$T_{1}$-proces}.

Místo RF pulsu larmorovského elektromagnetického
záření,\INDEX{záření!elektromagnetické} jenž anuluje podélnou magnetizaci,
můžeme použít delšího pulsu, jenž dokonce obrátí magnetizaci do opačného
směru, na $M_{z} = -M_{0}$, takže relaxace bude delší, ale charakteristický
čas $T_{1}$ bude stejný,
\begin{equation}
  M_{z}(t)=M_{0}\left(1-2\eto{-\frac{t}{T_{1}}}\right)\,. \label{spinlatb}
\end{equation}
Podélnou relaxaci ilustruje obrázek~\ref{longrelax} uprostřed.

Aplikujeme-li naopak kratší RF puls, jenž \uv{sklopí} magnetizaci jen do
roviny $xy$, bude vektor magnetizace konat precesi okolo osy $z$ s~Larmorovou
frekvencí. Jelikož pole $B_{0}$ nemůže být zcela exaktně homogenní, budou se
v~různých místech velmi mírně lišit i~odpovídající Larmorova
frekvence~(\ref{larmor}). Následkem toho dojde k~rozfázování jednotlivých
objemových elementů vytvářejících příčnou komponentu
$M_{\perp}$,\footnote{Názorně si lze tento proces představit jako rozevírání
  vějíře. Jednotlivá pera vějíře představují příspěvky do magnetizace
  objemového elementu. Původně zcela složený vějíř odpovídá příčné magnetizaci
  $M_{\perp0}$, po úplném rozevření do kruhu pak zcela relaxované, nulové
  příčné magnetizaci.} a postupnému zmenšování příčné komponenty podle rovnice
\begin{equation}
  M_{\perp}(t) = M_{\perp0}\,\eto{-\frac{t}{T_{2}}}\,,    \label{spinspin}
\end{equation}
kde $M_{\perp0}$ je počáteční hodnota příčné magnetizace a $T_{2}$ je
\emph{spin-spinový relaxační čas}, jenž je vždy menší nebo rovný $T_{1}$, tj.
příčná relaxace je rychlejší než podélná. Tento jev se nazývá
\emph{$T_{2}$-proces} a jeho časový průběh je ilustrován na
obrázku~\ref{longrelax} vpravo.  Vztahy (\ref{spinlat}) až (\ref{spinspin}) se
dají odvodit řešením \emph{Blochových rovnic}~\cite{Hor:LIVE:BasicsMRI:}.

Podrobný popis mechanismu \uv{sklápění} magnetizace pomocí
ra\-dio\-fre\-k\-venč\-ních pulsů poněkud vybočuje z~rámce této práce. Uveďme
pouze, že jej lze vysvětlit zavedením soustavy souřadnic $x'y'z'$ rotující
kolem osy $z'\equiv z$, jež je společná s~\uv{laboratorní} soustavou souřadnic
$xyz$, larmorovskou úhlovou rychlostí $\omega = 2\pi
f$~\cite{Eng:1995:POKMF:}.  Na tomto místě doporučujeme navštívit
stránku~\cite{Hor:LIVE:BasicsMRI:}, kde jsou velmi názorné animace ilustrující
relaxaci jaderných spinů, včetně popisu mechanismu \uv{sklápění} magnetizace.

Ve skutečnosti jsou oba procesy spojené a probíhají současně. Úhel otočení
vektoru magnetizace je úměrný energii RF pulsu. RF puls, jenž \uv{sklopí}
magnetizaci do roviny $xy$, se nazývá (ze zřejmých důvodů) $90^{\circ}$-puls,
a puls, jenž otočí magnetizaci do opačného směru, se nazývá
$180^{\circ}$-puls. Po odeznění pulsu, jenž provedl otočení nebo sklopení
magnetizace, dochází k~relaxaci zahrnující $T_{1}$-proces i $T_{2}$-proces.
Nedojde-li po aplikaci RF pulsu k~čerpání energie do spinů pomocí následného
RF pulsu, nazýváme volnou relaxaci \emph{FID (free induction decay)}.

Tvar RF pulsu je rovněž kritický a má vliv na přesnost výběru obrazového řezu.
Nejčastěji se používá zvonovitě tvarovaný puls obsahující 3 až 4 vlnové délky
$\lambda$ odpovídající Larmorově frekvenci~(\ref{larmor}), nebo puls tvarovaný
jako funkce $\sinc$ (sinus cardinalis). Takto tvarovaný RF puls má optimální
vlastnosti ve frekvenční doméně.  Tvar RF pulsu, jeho Fourierovo spektrum a
FID jsou zobrazeny na obrázku~\ref{mrifid}.

\begin{figure}[t]
\begin{center}

\includegraphics[width=.42\linewidth]{mrifid}
\end{center}
\caption[Ilustrace signálu FID]%
{\label{mrifid}Tvar $90^{\circ}$ RF pulsu, jeho PSD (výkonové Fourierovo
  spektrum) a volné doznívání příčné magnetizace (FID). Převzato
  z~\cite{Dow-Ken-Joh:2006:PhysDiagImag:}.}
\end{figure}

\begin{table}[b]
\caption[Relaxační doby pro různé tkáně]%
{\label{mrireltime}Relaxační doby $T_{1}$, $T_{2}$ pro různé tkáně. Relaxační
  doba $T_{1}$ závisí na magnetické indukci $B_{0}$ zhruba jako $B_{0}^{0.3}$
  a je zde zastoupena hodnotami pro tři typické používané intenzity
  $0.5\,\mathrm{T}$, $1.0\,\mathrm{T}$ a $1.5\,\mathrm{T}$.
Podle~\cite{Dow-Ken-Joh:2006:PhysDiagImag:}.}
\begin{center}
\small
\newcommand{\pad}{\hphantom{0}}
\begin{tabular}{lcccc}
\hline
Tkáň & \multicolumn{3}{c}{$T_{1}\,[\mathrm{ms}]$}
     & $T_{2}\,[\mathrm{ms}]$ \\ %\cline{2-4}
     & $0.5\,\mathrm{T}$ & $1.0\,\mathrm{T}$ & $1.5\,\mathrm{T}$ & \\\hline
Tuk  & 210 & 240 & 260 & 80 \\
Játra  & 350 & 420 & 500 & 40 \\
Ledviny  & 430 & 590 & 690 & 58 \\
Svaly  & 550 & 730 & 870 & 45 \\
Srdce  & 570 & 750 & 880 & 57 \\
Bílá hmota mozk. & 500 & 680 & 780 & 90 \\
Šedá hmota mozk. & 650 & 810 & 900 & 100\pad \\
Mozkomíšní mok  & 1800\pad & 2160\pad & 2400\pad & 160\pad \\
\hline
\end{tabular}
\end{center}
\end{table}

\begin{figure}[t]
\begin{minipage}{.48\linewidth}
\centering

\includegraphics[width=\linewidth]{ses}
\end{minipage}\hfill%
\begin{minipage}{.48\linewidth}
\centering

\includegraphics[width=\linewidth]{irs}
\end{minipage}
\caption[Sekvence RF pulsů \uv{spin-echo} a \uv{inversion-recovery}]%
{\label{mripseq}\emph{Vlevo:} Spin-echo sekvence.  Nejprve se aplikuje
  $90^{\circ}$ puls, jímž se magnetizace sklopí do roviny $xy$, dochází
  k~defázování a vysílání FID signálu. V~okamžiku, kdy je \uv{vějíř}
  magnetizací rozevřen právě do půlkruhu, se aplikuje $180^{\circ}$ puls, jímž
  se \uv{vějíř} překlopí okolo své hrany na opačnou stranu. Stejným
  mechanismem jako dochází k~defázování, nastává nyní sfázování (skládání
  \uv{vějíře}), jež vzápětí přejde opět k~defázování (rozkládání \uv{vějíře}),
  a~k~vyslání nejprve sílícího a poté slábnoucího echa. Časová odlehlost FID
  signálu a echa se označuje $T_{\mathrm{E}}$, a~celý proces se opakuje
  s~periodou $T_{\mathrm{R}}$. \emph{Vpravo:} Inversion recovery sekvence
  sestává z $180^{\circ}$ pulsu, jenž otočí magnetizaci do opačného směru,
  následovaného sekvencí spin-echo. Perioda $T_{\mathrm{I}}$ mezi
  $180^{\circ}$ pulsem a $90^{\circ}$ pulsem se nazývá \emph{inverzní čas}.
  Převzato z~\cite{Dow-Ken-Joh:2006:PhysDiagImag:}.}
\end{figure}

Podstatné je, že relaxační časy $T_{1}$ a $T_{2}$ jsou výsledkem vzájemného
působení rezonujících jader a jejich okolí a charakterizují chemické
vlastnosti a strukturu vyšetřované tkáně. Relaxační doby jsou proto často
výrazně odlišné pro různé typy tkáně a také pro zdravou a nádorovou tkáň.
Relaxační časy pro různé typy tkání sumarizuje tabulka~\ref{mrireltime}.

Při relaxaci dochází podle zákonů elektrodynamiky k~emisi elektromagnetického
záření,\INDEX{záření!elektromagnetické} jehož intenzita závisí na koncentraci
daného druhu rezonujících jader (v~našem případě protonů), na chemických
vlastnostech a struktuře tkáně.  Jsou-li RF pulsy při FID opakovány s~periodou
$T_{\mathrm{R}}$ dostatečně velikou, aby bylo zachováno volné doznívání, je
intenzita signálu emitovaného relaxujícími spiny úměrná jejich hustotě $\rho$
a v~čase doznívá jako $1-\eto{-T_{\mathrm{R}}/T_{1}}$.

V~praxi se používají pro zlepšení odstupu signálu od šumu jiné posloupnosti RF
pulsů. Nejběžnější jsou \emph{spin-echo} sekvence a \emph{inversion recovery}
sekvence, jež jsou vysvětleny na obrázku~\ref{mripseq}. Intenzity signálu
relaxujících spinů doznívají jako $(1-\eto{-T_{\mathrm{R}}/T_{1}})
\eto{-T_{\mathrm{E}}/T_{2}}$ pro sekvenci spin-echo a
$1-2\eto{-T_{\mathrm{I}}/T_{1}} +\eto{-T_{\mathrm{R}}/T_{1}}$ pro sekvenci
inversion recovery.  Významy časů $T_{\mathrm{E}}$, $T_{\mathrm{I}}$ a
$T_{\mathrm{R}}$ jsou vysvětleny v~popisce obrázku~\ref{mripseq}.

Záření emitované\INDEX{záření!emitované} relaxujícími spiny se zachycuje
v~přijímací cívce a analýzou získaného signálu lze rekonstruovat hustotu
rezonujících spinů uvnitř tkáně.  Jak se toho dosáhne je předmětem následující
sekce.

\subsection{Podstata MRI zobrazování}\label{mrimagbase}%///////////////////

V~medicínském MRI nás zajímá NMR hlavně signál z~vody a lipidů, komponentami
lidského těla obsahujícími nejvíce vodíku. Základním principem je rezonanční
vztah~(\ref{larmor}).  Pro názornost uvažujme na obrázku~\ref{mrihead}
znázorněnou lidskou hlavu obsahující tři oddělené malé oblasti s~vysokou
hustotou vodíkových atomů; jinde předpokládejme jejich výskyt nulový (ve
skutečnosti je vodík obsažen s~různou hustotou v~celém těle).

Pokud by naše fiktivní hlava byla vložena do homogenního magnetického pole
$B_{0}$, všechny tři oblasti by měly stejnou Larmorovu frekvenci a na NMR
spektru získaném Fourierovou analýzou NMR signálu relaxujících spinů by se
jejich rezonanční pík objevil při téže (Larmorově) frekvenci, jak je
znázorněno na obrázku~\ref{mrihead}. Situace se změní, jakmile homogenní
magnetické pole zaměníme \emph{gradientním magnetickým polem}, které na
obrázku roste zleva doprava.\footnote{Aplikovaný gradient magnetického pole by
  měl být tak velký, aby změna magnetické indukce $\Delta B$ v~rozsahu
  zobrazované oblasti odpovídala šířce pásma RF pulsů, která je okolo $\Delta
  f \sim 400\,\hz[k]$~\cite{Dow-Ken-Joh:2006:PhysDiagImag:}. Kdyby měl RF puls
  menší šířku pásma, bodům blíže k~okraji oblasti by se \uv{nedostávalo}
  záření o~rezonanční frekvenci. Zároveň je opět vidět, jak je důležitá
  znalost frekvenčního spektra RF pulsů, která určuje relativní citlivost při
  různých indukcích v~gradientním magnetickém poli.} Protože některé ze tří
oblasti s~vysokou hustotou vodíku leží v~magnetickém poli s~různou intenzitou,
dostaneme z~nich podle~(\ref{larmor}) NMR signál na různých frekvencích,
přičemž amplituda signálu je úměrná počtu vodíkových atomů. Tomuto postupu se
říká \emph{frekvenční kódování}. K~zobrazování to však pořád ještě nestačí,
protože dvě ze tří oblastí leží v~oblasti se stejnou magnetickou indukcí a ve
spektru je proto nelze rozlišit.

\begin{figure}[t]
\begin{minipage}{.48\linewidth}
\centering

\includegraphics[width=\linewidth]{headhomo}
\end{minipage}\hfill%
\begin{minipage}{.48\linewidth}
\centering

\includegraphics[width=\linewidth]{headgrad}
\end{minipage}
\caption[Frekvenční kódování NMR signálu]%
{\label{mrihead}Podstata frekvenčního kódování NMR signálu. \emph{Vlevo:}
  Fiktivní hlava se třemi oblastmi vysoké hustoty vodíkových atomů. Hlava se
  nachází v~homogenním magnetickém poli a ve všech třech oblastech je tedy
  stejná Larmorova frekvence. Frekvenční spektrum bude proto mít pík na jediné
  Larmorově frekvenci $B_{0}/\gamma$. \emph{Vlevo:} Přiložíme-li gradient
  magnetického pole, některé (ne nutně všechny) oblasti vysoké hustoty vodíku
  se ocitnou v~magnetickém poli s~různou indukcí a jejich Larmorova frekvence
  se liší. Spektrum již má více píků, jejichž amplituda je úměrná počtu
  vodíkových atomů. Upraveno podle~\cite{Hor:LIVE:BasicsMRI:}.}
\end{figure}

Naznačený problém řeší metoda \emph{zpětné projekce (backprojection imaging)},
historicky první MRI metoda zavedená
Lauterburem~\cite{Lau:1973:NATURE:}. Metoda zpětné projekce je rozšířením
frekvenčního kódování v~tom smyslu, že směr gradientu magnetického pole se
postupně natáčí s~určitým krokem, a~pro každý krok se změří spektrum NMR
signálu.

Rovinu řezu zafixujeme dalším, silnějším gradientním magnetickým
polem.\footnote{Silnějším proto, aby se oblast rezonance $\Delta B = \Delta
  f/\gamma$ zúžila na pokud možno tenkou rovinu.}  Jestliže na
obrázku~\ref{mrihead} leží souřadná rovina $xy$ v~rovině kresby a osa $z$ je
k~ní kolmá, a jestliže má gradientní magnetické pole nenulovou pouze
$z$-složku, můžeme ji zapsat ve tvaru $B_{z} = B_{0} + z G_{z}$, kde $B_{0}$
je indukce v~určitém referenčním bodě, \emph{isocentru} magnetu, a konstantní
veličina $G_{z}$ reprezentuje rychlost změny (gradient) magnetického pole.
Rezonance nastává podle~(\ref{larmor}) tam, kde je velikost $B_{z}$ magnetické
indukce rovna Larmorově frekvenci protonu dělené gyromagnetickým poměrem
$\gamma$, takže pro geometrický útvar $B_{z}=f/\gamma = \mathrm{const}$,
v~němž dochází k~rezonanci, máme
\begin{equation}
  zG_{z} = \frac{f}{\gamma} - B_{0}
    = \mathrm{const}\,.                                   \label{gradfield}
\end{equation}
To je rovnice roviny kolmé k~ose $z$, tj. sagitálního řezu hlavou. Gradientní
pole $\vec{G}$ se dá obecně nastavovat složitějším způsobem, než jsme právě
popsali, a tím měnit sklon a polohu roviny řezu.\footnote{Není principiálního
  důvodu k~tomu, aby geometrickým útvarem byla nerovinná
  plocha~\cite{Mad:2007:SiemensMRI:OralCom}. Autorce se nepodařilo zjistit,
  zda se nerovinných řezů v~praxi užívá.} Ve skutečnosti dochází k~rezonanci
v~tenké rovinné vrstvě vlivem tepelného pohybu molekul a vlivem rozšíření
spektrální čáry RF pulsů.

\begin{figure}[t]
\begin{minipage}{.48\linewidth}
\centering

\includegraphics[width=\linewidth]{bkproja}
\end{minipage}\hfill%
\begin{minipage}{.48\linewidth}
\centering

\includegraphics[width=\linewidth]{bkprojb}
\end{minipage}
\caption[Rekonstrukce obrazu pomocí techniky zpětné projekce]%
{\label{mribkproj}Podstata rekonstrukce obrazu pomocí techniky zpětné
  projekce. \emph{Vlevo:} Pomocí gradientního pole (\ref{gradfield}) ve směru
  kolmém k~rovině obrázku se zafixuje rovina řezu (zde sagitální), a pomocí
  gradientů v~rovině obrázku natáčených s~určitým krokem (zde pro názornost
  $45^{\circ}$) od $0^{\circ}$ do $360^{\circ}$ se v~každém kroku sejme NMR
  spektrum.  \emph{Vpravo:} Ze znalosti všech spekter lze pomocí počítačově
  implementované Radonovy transformace zrekonstruovat mapu sagitálního řezu
  vybraného pomocí gradientu kolmého k~rovině obrázku.  Upraveno
  podle~\cite{Hor:LIVE:BasicsMRI:}.}
\end{figure}

Odtud mimochodem plyne důležitost tvarování RF pulsů. Kdybychom tyto pulsy
vytvarovali s~obdélníkovou obálkou, jejich frekvenční spektrum by mělo tvar
funkce $\sinc$ se známými oscilačními laloky. Frekvence by měla velký rozptyl
a v~důsledku toho by byl velký rozptyl i~lokalizaci míst, kde dochází
k~rezonanci; jinými slovy, rovina řezu by měla větší tloušťku, což by mohlo
vést ke ztrátě ostrosti a detailu. Jestliže však RF pulsy příhodně
vytvarujeme, například tak, aby jejich obálka měla tvar funkce $\sinc$, bude
mít naopak frekvenční spektrum pulsů obdélníkový tvar a rovina zobrazení bude
dobře lokalizovaná. RF pulsy s~gaussovským profilem poskytnou opět gaussovské
frekvenční spektrum, tedy dobře lokalizovanou rovinu řezu. Tvarování RF pulsů
tak lze využít k~regulaci síly řezu.

Ze znalosti NMR spekter při všech krocích pak lze numericky implementovanou
inverzní Radonovou transformací (zpětná projekce), podobně jako u~počítačové
tomografie, pomocí počítače zrekonstruovat vybraný řez, jak je znázorněno na
obrázcích~\ref{mribkproj} a~\ref{mriloc}.

Ačkoli je technika frekvenčního kódování a zpětné projekce poměrně jednoduchá
a názorná, v~současných MRI modalitách se nepoužívá. Místo ní se aplikuje
další kategorie zobrazování nazývaná \emph{dvourozměrné fourierovské
  zobrazování} objevené \designation{R. Ernstem} a založené na Fourierově
transformaci. Podrobný výklad této poměrně složité techniky se vymyká rozsahu
bakalářské práce a odkazujeme proto např.  na~\cite{Hor:LIVE:BasicsMRI:}. Zde
uvedeme jen základní ideu~\cite{Hau:1990:CZEJPA:}.

\begin{figure}[t]
\begin{minipage}{.48\linewidth}
\centering

\includegraphics[width=\linewidth]{axhead1}
\end{minipage}\hfill%
\begin{minipage}{.48\linewidth}
\centering

\includegraphics[width=\linewidth]{axhead2}
\end{minipage}
\caption[MRI axiální zobrazení lidského mozku s~patologií]%
{\label{mriaxhead}Axiální (N)MRI zobrazení lidského mozku s~patologií.
  \emph{Vlevo:} $T_{\mathrm{R}}/T_{\mathrm{E}} = 5500/105\,\mathrm{ms}$, 2~Nex
  (počet zprůměrovaných obrazů), matice $512\times256$. \emph{Vpravo:}
  $T_{\mathrm{R}}/T_{\mathrm{E}} = 450/14\,\mathrm{ms}$, 1~Nex, matice
  $256\times192$.  Převzato z~\cite{Hor:LIVE:BasicsMRI:}.}
\end{figure}

Představme si, že na obrázku~\ref{mriloc} je osa $z$ orientována vertikálně
(rovnoběžně se hřbetem vazby), osa $y$ leží v~rovině obrázku a osa $x$ kolmo
k~obrázku. \uv{Výšku} transverzálního řezu, tj. $z$-ovou polohu řezu,
regulujeme stálým gradientem $G_{z}$. Do hry nyní vstupují přechodné gradienty
ve zbývajících osách $x$ a $y$. Nejprve je aplikován gradient $G_{x}$ ve směru
osy $x$ po dobu $t_{x}$, potom gradient $G_{y}$ ve směru osy $y$ po dobu
$t_{y}$. Označme hustotu vodíkových jader ve vybraném řezu (pomocí gradientu
$G_{z}$) $\rho(x,y)$. Jsou-li časy $t_{x,y}$ přikládání gradientů $G_{x,y}$
malé ve srovnání s~relaxačními časy, pak je časový průběh signálu NMR
obrazového elementu $\dif x\dif y$ v~bodě $x,y$ úměrný
\begin{equation}
  \rho\,\eto{\iu\gamma Bt}\,\dif x\dif y\,.            \label{mrisigtevol}
\end{equation}
V~časovém intervalu $t_{x}$ je magnetická indukce $B=B_{0}+xG_{x}$,
v~následující době $t_{y}$ platí $B=B_{0}+yG_{y}$, kde $B_{0}$ označuje
hodnotu v~místě $x=y=0$. Dosadíme-li oba poslední výrazy do
vztahu~(\ref{mrisigtevol}) a členy obsahující na souřadnicích nezávislý
konstantní člen $B_{0}$ zahrneme do konstanty úměrnosti, dostaneme pro
velikost signálu NMR hodnotu úměrnou $\rho\,\eto{\iu\gamma
  (xG_{x}t_{x}+yG_{y}t_{y})}\,\dif x\dif y$, a pro celý řez
\begin{equation}
  s(t_{x},t_{y})
    = \iint\rho(x,y)\,\eto{\iu\gamma (xG_{x}t_{x}+yG_{y}t_{y})}\,
    \dif x\dif y\,.                                     \label{mrisigslice}
\end{equation}
Zavedeme-li substituce $\omega_{x}=\gamma xG_{x}$, $\omega_{y}=\gamma yG_{y}$
a $A=1/\gamma^{2}G_{x}G_{y}$, můžeme vztah~(\ref{mrisigslice}) přepsat jako
\begin{equation}
  s(t_{x},t_{y})
    = A\iint\rho(\omega_{x},\omega_{y})\,
    \eto{\iu(\omega_{x}t_{x}+\omega_{y}t_{y})}\,
    \dif\omega_{x}\dif\omega_{y}\,.                           \label{mrift}
\end{equation}
Ale to je dvourozměrná Fourierova transformace mezi měřitelnou veličinou
$s(t_{x},t_{y})$ a hledaným rozložením vodíkových jader v~řezu
$\rho(\omega_{x},\omega_{y})$, potažmo $\rho(x,y)$.\footnote{Frekvence spínání
  gradientních cívek, jež leží ve slyšitelné akustické oblasti, zapříčiňuje
  \uv{kulometnou palbu}, obvyklý zvukový doprovod MRI vyšetření.} Tato metoda
ukazuje, že Fourierova transformace nevystupuje vždy jen v~roli \uv{pomocného}
nebo dodatečného zpracování, jako jsou filtrace a různé úpravy signálu, ale
může být nativní, neoddělitelnou součástí převodu primárního signálu
$s(t_{x},t_{y})$ na obrazový výstup $\rho(x,y)$.

Obraz získaný metodou MRI je mapou protonové hustoty ve vybraném řezu.
Přesněji, na MRI zobrazeních jsou zachyceny dvě základní informace:
\begin{itemize}
\item Distribuce hustoty jader vytvářejících NMR (obvykle vodíkových jader).
\item Distribuce relaxačních časů souvisejících s~chemickým složením
  a strukturním stavem tkáně v~jednotlivých místech.
\end{itemize}
Numerickou analýzou relaxačních časů NMR signálu lze vytvářet mapy
v~relaxačních časech $T_{1}$ a $T_{2}$ (označovaných jako $T_{1}$ či
$T_{2}$-vážené obrazy).  Soustava obrazů různých rovnoběžných řezů pak vytváří
3D~obraz vyšetřované oblasti, z~něhož lze pomocí metod počítačové
grafiky~\cite{Zar-etal:2004:ModPocGraf:} vytvářet obrazy libovolných řezů.
Oba výše uvedené faktory jsou jiné nejen pro různé druhy tkání (viz
tabulka~\ref{mrireltime}), ale liší se i~v~závislosti na fyziologickém či
patologickém stavu téže tkáně. To vše spolu se skutečností absence
ionizujícího záření\INDEX{záření!ionizující} činí z~MRI významnou
diagnostickou metodu. Na obrázku~\ref{mriaxhead} je příklad axiálního řezu
lidskou hlavou s~patologií v~mozku; tloušťka řezu $5\,\mathrm{mm}$
$\mathrm{FOV}=22\,\mathrm{cm}$ (field of view~\cite{Hor:LIVE:BasicsMRI:}).

\endinput

%% 
%% End of file `x_diag.tex'.
